\begin{abstract}
Frequent co-purchases describe historical baskets, but their persistence in later transactions is an empirical question. We study this gap using UCI Online Retail with separate discovery, validation and test periods. A validation-selected cohort of \SelectedRules{} rules achieves \MeanRetention{} mean monthly retention in four held-out months; \AllMonthsRules{} rules satisfy every monthly requirement throughout. The interpretation depends on observability: low-support cohorts frequently lack adequate monthly evidence, while restricting evaluation to orders with recorded customer IDs lowers retention to \KnownRetention{}. Identical bounded Apriori and FP-Growth searches return the same itemsets and support counts at modest local computational cost. The study contributes an auditable temporal evaluation and diagnoses how thresholds, observation coverage and order composition affect apparent stability. It does not establish causal retail utility or statistical significance of individual rules.
\end{abstract}

\section{Introduction}
Association rules summarize products that occur together in a shopping basket. Such patterns can inform candidate bundles or merchandising decisions, but a rule with high historical confidence may weaken when the order mix, product availability or purchasing context changes. The practical question is whether an association remains sufficiently common and reliable after discovery, rather than whether it fits the period in which it was found.

\textbf{Research questions.} We ask (RQ1) how often a fixed set of discovered rules satisfies the same requirements in later months; (RQ2) how support and confidence thresholds trade rule quantity against retention and adequate evidence; and (RQ3) how two established mining implementations compare under identical inputs and search bounds. UCI Online Retail~\citep{chen2015} supplies a historical, invoice-level case study. Rules are discovered in December 2010--May 2011, settings are selected on June--July, and the frozen cohort is evaluated on August--November.

\textbf{Relation to prior work.} Apriori and FP-Growth address efficient frequent-itemset enumeration through different search structures~\citep{agrawal1994,han2000}. We use them as equivalent discovery engines, not as competing predictors. Temporal evaluation is motivated by the broader problem of changing distributions~\citep{gama2014}. Separately, statistically sound pattern discovery studies false discoveries and multiplicity~\citep{hamalainen2019}; satisfying support, confidence and lift thresholds is not a substitute for those tests. Our focus is operational persistence under explicit criteria, rather than a drift detector or a new significance procedure.

\textbf{Contributions and scope.} First, we provide a reproducible chronology that separates discovery, selection and evaluation. Second, we distinguish the availability of evidence from success under metric thresholds, exposing a major limitation of low-support cohorts. Third, fixed-cohort population checks and an exact algorithm-equivalence benchmark make the interpretation auditable. The evidence concerns one retailer and four test months; generalization and commercial utility require further evaluation.

All comparisons use the same cleaned baskets and a common monthly definition of success. Changing the cohort or evaluation criteria between periods could make a comparison appear more stable without preserving the same relationships. We therefore retain unsuccessful and unobservable rules in every test denominator.

